% Context from chapters-tex/approximation.tex; for mathematical reference, not compilation.
\section{Signal and Image Modeling}
\label{sec-signal-models}

A signal model specifies a constraint $f \in \Theta$, where $\Theta \subset \Ldeux([0,1]^d)$ is the class of signals of interest.
Figure \ref{fig-approximation-class} shows different image models, which we describe below.

\myfigure{
\tabquatre{
\image{approximation}{.24}{approximation-class-regular}&
\image{approximation}{.24}{approximation-class-bv}&
\image{approximation}{.24}{approximation-class-cartoon}&
\image{approximation}{.24}{approximation-class-natural}\\
Regular & Bounded variation & Cartoon & Natural
}
                                                            
}{ 
	Examples of image models.  
}{fig-approximation-class}



                                           
\subsection{Uniformly Smooth Signals and Images}
\label{subsec-smooth-class}

  
\paragraph{Signals with derivatives.}

The simplest model consists of uniformly smooth signals with bounded derivatives
\eql{\label{eq-smooth-model} 
	\Th = \enscond{ f \in \Ldeux([0,1]^d) }{ \norm{f}_{\Cal} \leq C },
}
where $C>0$ is fixed. For integer $\alpha\geq1$, in one dimension
\eq{
	\norm{f}_{\Cal} \eqdef \umax{0\leq k\leq\al} \normb{ \frac{\d^k f}{\d t^k} }_{\infty}.
}
In higher dimensions, take the maximum over all partial derivatives of total order at most $\alpha$. For noninteger $\alpha=m+\beta$, $0<\beta<1$, use the usual H\"older norm: derivatives through order $m$ are bounded, and order-$m$ derivatives are $\beta$-H\"older continuous.

  
\paragraph{Sobolev smooth signals and images.}

For integer $\alpha\geq1$, a $\Cal$ signal in the model~\eqref{eq-smooth-model} has derivatives with bounded energy:
\eq{
	\frac{\d^\al f}{\d t^\al}(t) = f^{(\al)}(t) \in \Ldeux([0,1]). 
}	
For integer $\alpha$ and periodic boundary conditions, use the identity
\eq{
	\hat f^{(\al)}_m = (2i\pi m)^\al \hat f_m
}
Here $\hat f$ denotes Fourier coefficients as in~\eqref{eq-defn-fourier-coeffs}, now on $\RR/\ZZ$ rather than $\RR/2\pi\ZZ$:
\eq{
	\hat f_n \eqdef \int_0^1 e^{-2\imath\pi n x} f(x) \d x, 
}
This motivates the homogeneous Sobolev seminorm
\eql{\label{eq-sobolev-norm}
	\norm{f}_{\text{Sob}(\al)}^2 = \sum_{m \in \ZZ} |2\pi m|^{2\al} |\dotp{f}{e_m}|^2,
}
so that $\norm{f}_{\text{Sob}(\al)} = \norm{f^{(\al)}}$ for smooth periodic functions. The same seminorm applies when the derivatives exist only in the distributional sense and belong to $\Ldeux([0,1])$.

This definition extends to distributions and signals $f \in \Ldeux([0,1]^d)$ of arbitrary dimension $d$ as
\eql{\label{eq-sobolev-norm-highdim}
	\norm{f}_{\text{Sob}(\al)}^2 = \sum_{m \in \ZZ^d} (2\pi \norm{m})^{2\al} |\dotp{f}{e_m}|^2,
}

The periodic Sobolev model
\eql{\label{eq-sobolev-